\documentclass{report}
\usepackage{amsmath,amssymb}
\setlength{\parindent}{0mm}
\setlength{\parskip}{1em}
\begin{document}
\begin{center}
\rule{6in}{1pt} \
{\large John Shadid \\
{\bf Scalable Solution of Continuum Plasma Physics Models by Approximate Block Factorization Preconditioners}}

Sandia National Laboratories \\ Computational Mathematics Department \\ Sandia National Laboratories \\ PO Box 5800 MS 1321 \\ Albuquerque \\ NM 87185 � 1321
\\
{\tt jnshadi@sandia.gov}\\
Edward Phillips\\
Eric Cyr\\
Roger Pawlowski\end{center}

The mathematical basis for the continuum modeling of plasma physics systems
is the solution of the governing partial differential equations (PDEs)
describing conservation of mass, momentum, and energy, along with various
forms of approximations to Maxwell's equations. This PDE system is
non-self adjoint, strongly-coupled, highly-nonlinear, and characterized
by physical phenomena that span a very large range of length- and
time-scales. To enable accurate and stable approximation of these systems
a range of spatial and temporal discretization methods are commonly
employed. In the context of finite element spatial discretization methods
these include mixed integration, stabilized methods and
structure-preserving (physics compatible) approaches. For effective time
integration for these systems some form of implicitness is required. Two
well-structured approaches, of recent interest, are fully-implicit and
implicit-explicit (IMEX) type methods. Clearly the desire to accommodate
disparate spatial discretizations, and allow the flexible assignment of
mechanisms as explicit or implicit operators, implies a wide variation in
unknown coupling, ordering, the nonzero block structure, and the
conditioning of the implicit system. These characteristics make the
scalable and efficient iterative solution of these systems extremely
challenging.

This talk considers the development and evaluation of approaches based on
approximate block factorization (ABF) and physics-based preconditioning
approaches. These methods reduce the coupled systems into a set of
simplified systems to which multilevel methods are applied. A critical
aspect of these methods is the development of approximate Schur
complement operators that encode the critical cross-coupling physics of
the system [1-3]. To demonstrate the flexibility and performance of these
methods we consider application of these techniques to resistive MHD and
multiphysics MHD models for challenging prototype systems.
In this context robustness, efficiency, and the parallel and algorithmic
scaling of the preconditioning methods are discussed.

[1] E. C. Cyr, J. N. Shadid, R. S. Tuminaro, R. P. Pawlowski, and L.
Chacon, A new approximate block factorization preconditioner for 2D
incompressible (reduced) resistive mhd, SIAM Journal on Scientific
Computing, 35:B701-B730, 2013

[2] E. G. Phillips, H. C. Elman, E. C. Cyr, J. N. Shadid, R.P. Pawlowski,
A Block Preconditioner for an Exact Penalty Formulation for Stationary
MHD, SIAM J. Sci. Comput. 36-6 (2014), pp. B930-B951

[3] E. C. Cyr, J. N. Shadid, and R. S. Tuminaro, �Teko an abstract block
preconditioning capability with concrete example applications to
Navier-Stokes and resistive MHD, Accepted in SISC

*This work was supported by the DOE Office of Science Advanced Scientific
Computing Research - Applied Math Research program at Sandia National
Laboratories.


\end{document}
